\documentclass[11pt]{article}
\usepackage[utf8]{inputenc}
\usepackage[T1]{fontenc}
\usepackage[english,provide=*]{babel}
\usepackage[a4paper,margin=2.6cm]{geometry}
\usepackage{lmodern}
\usepackage{microtype}
\usepackage{amsmath,amssymb,amsfonts,mathtools}
\usepackage{booktabs,tabularx,array}
\usepackage{enumitem}
\usepackage{float}
\usepackage{graphicx}
\usepackage{xcolor}
\usepackage[colorlinks=true,linkcolor=blue,urlcolor=blue,citecolor=blue]{hyperref}
\hypersetup{
  pdftitle={Online Resource 1 - Supplementary Information for Extrinsic Newton Method with Retractions for Zeros of Vector Fields on Regular Level-Set Submanifolds},
  pdfauthor={Rodrigo Alberto Castro Marin and Leslie San Martin},
  pdfsubject={Supplementary Information for Calcolo}
}
\usepackage[capitalize,nameinlink]{cleveref}
\crefname{table}{table}{tables}
\Crefname{table}{Table}{Tables}
\crefname{equation}{equation}{equations}
\Crefname{equation}{Equation}{Equations}
\newcommand{\R}{\mathbb{R}}
\newcommand{\M}{\mathcal{M}}
\newcommand{\T}{\mathrm{T}}
\newcommand{\D}{\mathrm{D}}
\newcommand{\Proj}{\mathrm{P}}
\newcommand{\Exp}{\mathrm{Exp}}
\newcommand{\Retr}{\mathrm{R}}
\newcommand{\PT}{\mathcal{P}}
\DeclareMathOperator{\sym}{sym}
\renewcommand{\thesection}{S\arabic{section}}
\renewcommand{\thesubsection}{S\arabic{section}.\arabic{subsection}}
\renewcommand{\theequation}{S\arabic{equation}}
\renewcommand{\thetable}{S\arabic{table}}
\renewcommand{\thefigure}{S\arabic{figure}}
\title{Online Resource 1\\
Supplementary Information for\\
\textit{Extrinsic Newton Method with Retractions\\
for Zeros of Vector Fields on Regular Level-Set Submanifolds}}
\author{Rodrigo Alberto Castro Mar\'in\thanks{Corresponding author. Email: \href{mailto:rodrigo.castrom@uv.cl}{rodrigo.castrom@uv.cl}.}
\and Leslie San Mart\'in\thanks{Email: \href{mailto:l.sanmartindonoso1@uandresbello.edu}{l.sanmartindonoso1@uandresbello.edu}.}}
\date{}
\begin{document}
\maketitle
\begin{center}
\small
\textbf{Journal:} \textit{Calcolo}\\
Institute of Mathematics, Universidad de Valpara\'iso, Valpara\'iso, Chile
\end{center}

\section*{Contents of Online Resource 1}
This online resource complements the main article. It presents an example on
the product torus $S^1\times S^1$ and a nonlinear field on
$\operatorname{St}(3,2)$. It also contains the complete trajectories of
residuals and cumulative times for the eight experiments, the moderate-precision
timing protocol, the controls for the Riemannian exponential on
$SL(2,\R)$, and the information required to reproduce the computations.

In all comparisons, the same equation is used to compute the Newton
direction, evaluated at the iterate of each variant. The differences arise
exclusively from the update by means of the exponential map or the retraction
under consideration.

\section[Example on the product torus]{Example on the product torus $S^1\times S^1$}
\label{sec:product_torus}
% ============================================================

Consider the product torus
\begin{equation}\label{eq:product_torus}
    \M
    =
    S^1\times S^1
    =
    \left\{
        (x,y,z,w)\in\R^4:
        x^2+y^2=1,\quad z^2+w^2=1
    \right\}.
\end{equation}
The two constraints are grouped in the mapping
\begin{equation}\label{eq:F_product_torus}
    F(x,y,z,w)
    =
    \begin{pmatrix}
        x^2+y^2-1\\
        z^2+w^2-1
    \end{pmatrix}.
\end{equation}
Its derivative is
\begin{equation}\label{eq:DF_product_torus}
    \D F(p)
    =
    \begin{pmatrix}
        2x&2y&0&0\\
        0&0&2z&2w
    \end{pmatrix}.
\end{equation}
Since $(x,y)\neq(0,0)$ and $(z,w)\neq(0,0)$ on $\M$, the two rows of
$\D F(p)$ are linearly independent. Therefore, $0$ is a regular value and
$\M$ is a regular submanifold of dimension two.

% ------------------------------------------------------------
\subsection*{Tangent space, pseudoinverse, and projector}
% ------------------------------------------------------------

The tangent space is
\begin{equation}\label{eq:tangent_product_torus}
    \T_p\M
    =
    \left\{
        (u_1,u_2,u_3,u_4)^\top:
        xu_1+yu_2=0,\quad
        zu_3+wu_4=0
    \right\}.
\end{equation}
Since
\[
    \D F(p)\D F(p)^\top=4I_2,
\]
the Moore--Penrose pseudoinverse is
\begin{equation}\label{eq:pseudoinverse_product_torus}
    \D F(p)^\dagger
    =
    \frac12
    \begin{pmatrix}
        x&0\\
        y&0\\
        0&z\\
        0&w
    \end{pmatrix}.
\end{equation}
Therefore, the orthogonal projector onto $\T_p\M$ is
\begin{equation}\label{eq:projector_product_torus}
    \Proj_p
    =
    I_4-\D F(p)^\dagger\D F(p)
    =
    \begin{pmatrix}
        y^2&-xy&0&0\\
        -xy&x^2&0&0\\
        0&0&w^2&-zw\\
        0&0&-zw&z^2
    \end{pmatrix}.
\end{equation}

A natural orthonormal basis of $\T_p\M$ is
\begin{equation}\label{eq:basis_product_torus}
    e_\theta(p)=(-y,x,0,0)^\top,
    \qquad
    e_\varphi(p)=(0,0,-w,z)^\top.
\end{equation}
We write
\begin{equation}\label{eq:Z_product_torus}
    Z_p
    =
    \begin{pmatrix}
        -y&0\\
        x&0\\
        0&-w\\
        0&z
    \end{pmatrix}.
\end{equation}
Then
\[
    Z_p^\top Z_p=I_2,
    \qquad
    \Proj_p=Z_pZ_p^\top.
\]

% ------------------------------------------------------------
\subsection*{Definition of the field}
% ------------------------------------------------------------

We define
\begin{align}
    f_1(y,w)
    &=
    y+\frac12w+y^2+\frac15yw,
    \label{eq:f1_product_torus}\\
    f_2(y,w)
    &=
    -\frac13y+w+w^2+\frac14yw.
    \label{eq:f2_product_torus}
\end{align}
We consider the ambient field
\begin{equation}\label{eq:field_product_torus}
    \widetilde X(x,y,z,w)
    =
    f_1(y,w)
    \begin{pmatrix}
        -y\\x\\0\\0
    \end{pmatrix}
    +
    f_2(y,w)
    \begin{pmatrix}
        0\\0\\-w\\z
    \end{pmatrix}.
\end{equation}
In coordinates,
\begin{equation}\label{eq:field_product_torus_components}
    \widetilde X(x,y,z,w)
    =
    \begin{pmatrix}
        -yf_1(y,w)\\
        xf_1(y,w)\\
        -wf_2(y,w)\\
        zf_2(y,w)
    \end{pmatrix}.
\end{equation}

Let us verify tangency. By \eqref{eq:DF_product_torus},
\[
\begin{aligned}
    \D F(p)\widetilde X(p)
    &=
    \begin{pmatrix}
        2x&2y&0&0\\
        0&0&2z&2w
    \end{pmatrix}
    \begin{pmatrix}
        -yf_1\\xf_1\\-wf_2\\zf_2
    \end{pmatrix}\\
    &=
    \begin{pmatrix}
        -2xyf_1+2xyf_1\\
        -2zwf_2+2zwf_2
    \end{pmatrix}\\
    &=0.
\end{aligned}
\]
Therefore,
\[
    X:=\widetilde X|_\M
\]
is a tangent vector field on $\M$.

The point
\begin{equation}\label{eq:zero_product_torus}
    p_\ast=(1,0,1,0)^\top
\end{equation}
is a zero, since
\[
    f_1(0,0)=f_2(0,0)=0.
\]

% ------------------------------------------------------------
\subsection*{Ambient Jacobian}
% ------------------------------------------------------------

The derivatives of the coefficients are
\begin{align*}
    \partial_yf_1
    &=
    1+2y+\frac15w,
    &
    \partial_wf_1
    &=
    \frac12+\frac15y,\\
    \partial_yf_2
    &=
    -\frac13+\frac14w,
    &
    \partial_wf_2
    &=
    1+2w+\frac14y.
\end{align*}
For brevity, we write
\[
    f_{1y}=\partial_yf_1,\quad
    f_{1w}=\partial_wf_1,\quad
    f_{2y}=\partial_yf_2,\quad
    f_{2w}=\partial_wf_2.
\]
The ambient Jacobian of \eqref{eq:field_product_torus_components} is
\begin{equation}\label{eq:jacobian_product_torus}
    \D\widetilde X(p)
    =
    \begin{pmatrix}
        0&-f_1-yf_{1y}&0&-yf_{1w}\\
        f_1&xf_{1y}&0&xf_{1w}\\
        0&-wf_{2y}&0&-f_2-wf_{2w}\\
        0&zf_{2y}&f_2&zf_{2w}
    \end{pmatrix}.
\end{equation}

The extrinsic covariant derivative is
\[
    \nabla_\eta X(p)
    =
    \Proj_p\D\widetilde X(p)(\eta).
\]

% ------------------------------------------------------------
\subsection*{Reduced Newton system}
% ------------------------------------------------------------

Let
\[
    \eta
    =
    \alpha e_\theta(p)+\beta e_\varphi(p)
    =
    Z_p
    \begin{pmatrix}
        \alpha\\\beta
    \end{pmatrix}.
\]
The reduced matrix is
\[
    H(p)
    =
    Z_p^\top
    \Proj_p\D\widetilde X(p)Z_p.
\]
A direct computation gives
\begin{equation}\label{eq:H_product_torus}
    H(p)
    =
    \begin{pmatrix}
        x\left(1+2y+\frac15w\right)
        &
        z\left(\frac12+\frac15y\right)
        \\[2mm]
        x\left(-\frac13+\frac14w\right)
        &
        z\left(1+2w+\frac14y\right)
    \end{pmatrix}.
\end{equation}
Likewise,
\begin{equation}\label{eq:reduced_residual_product_torus}
    Z_p^\top X(p)
    =
    \begin{pmatrix}
        f_1(y,w)\\
        f_2(y,w)
    \end{pmatrix}.
\end{equation}

The Newton direction is obtained by solving the reduced system
\begin{equation}\label{eq:newton_system_product_torus}
    H(p)
    \begin{pmatrix}
        \alpha\\\beta
    \end{pmatrix}
    =
    -
    \begin{pmatrix}
        f_1(y,w)\\
        f_2(y,w)
    \end{pmatrix}.
\end{equation}
If
\[
    \Delta(p)=\det H(p)\neq0,
\]
then
\begin{align}
    \alpha
    &=
    \frac{
        -h_{22}f_1+h_{12}f_2
    }{\Delta(p)},
    \label{eq:alpha_product_torus}\\
    \beta
    &=
    \frac{
        h_{21}f_1-h_{11}f_2
    }{\Delta(p)}.
    \label{eq:beta_product_torus}
\end{align}
At the zero $p_\ast$,
\begin{equation}\label{eq:H_star_product_torus}
    H(p_\ast)
    =
    \begin{pmatrix}
        1&\frac12\\[1mm]
        -\frac13&1
    \end{pmatrix},
\end{equation}
and
\begin{equation}\label{eq:det_H_star_product_torus}
    \det H(p_\ast)
    =
    1+\frac16
    =
    \frac76
    \neq0.
\end{equation}
Therefore, $p_\ast$ is a regular zero. The inverse function theorem
guarantees that it is the unique zero of the field in a sufficiently small
neighborhood of $p_\ast$.

% ------------------------------------------------------------
\subsection*{Angular interpretation}
% ------------------------------------------------------------

We can parametrize the torus by
\begin{equation}\label{eq:parametrization_product_torus}
    p(\theta,\varphi)
    =
    \bigl(
        \cos\theta,\sin\theta,
        \cos\varphi,\sin\varphi
    \bigr).
\end{equation}
Then
\[
    e_\theta=\frac{\partial p}{\partial\theta},
    \qquad
    e_\varphi=\frac{\partial p}{\partial\varphi}.
\]
The field is written as
\[
    X(p(\theta,\varphi))
    =
    f_1(\sin\theta,\sin\varphi)e_\theta
    +
    f_2(\sin\theta,\sin\varphi)e_\varphi.
\]
The matrix \eqref{eq:H_product_torus} is precisely the Jacobian
\[
    \frac{\partial(f_1,f_2)}
         {\partial(\theta,\varphi)}.
\]
This equality is used in the code only as a verification. The
direction is computed with the extrinsic formula
\[
    Z_p^\top \Proj_p\D\widetilde X(p)Z_p.
\]

% ------------------------------------------------------------
\subsection*{Exponential map and retraction}
% ------------------------------------------------------------

Since the torus is the product of two unit circles, the exponential map
acts by a rotation in each factor. For
\[
    \eta=\alpha e_\theta+\beta e_\varphi,
\]
we have
\begin{equation}\label{eq:exp_product_torus}
\begin{aligned}
    \Exp_p(\eta)
    =
    \bigl(
        &x\cos\alpha-y\sin\alpha,\,
        y\cos\alpha+x\sin\alpha,\\
        &z\cos\beta-w\sin\beta,\,
        w\cos\beta+z\sin\beta
    \bigr).
\end{aligned}
\end{equation}

A closed-form retraction is obtained by normalizing each pair separately:
\begin{equation}\label{eq:retraction_product_torus}
\begin{aligned}
    R_p(\eta)
    =
    \left(
        \frac{x-\alpha y}{\sqrt{1+\alpha^2}},
        \frac{y+\alpha x}{\sqrt{1+\alpha^2}},
        \frac{z-\beta w}{\sqrt{1+\beta^2}},
        \frac{w+\beta z}{\sqrt{1+\beta^2}}
    \right).
\end{aligned}
\end{equation}
This formula applies the blockwise homogeneous retraction to the two
circular constraints. In particular,
\[
    R_p(0)=p,
    \qquad
    \D R_p(0)=I_{\T_p\M}.
\]


% ------------------------------------------------------------

\subsection*{Updated numerical results}
\begin{table}[H]
\centering
\scriptsize
\setlength{\tabcolsep}{4.2pt}
\begin{tabular}{rccrr}
\toprule
$k$ & $\|X(p_k)\|_{\Exp}$ & $\|X(p_k)\|_R$
& $t_{\Exp,k}$ (s) & $t_{R,k}$ (s)\\
\midrule
0 & $5.87335\times10^{-1}$ & $5.87335\times10^{-1}$ & 0.000439 & 0.000431 \\
1 & $3.135593\times10^{-1}$ & $2.353352\times10^{-1}$ & 0.012963 & 0.005805 \\
2 & $3.721479\times10^{-2}$ & $2.590881\times10^{-2}$ & 0.024704 & 0.010646 \\
3 & $7.835498\times10^{-4}$ & $4.2763\times10^{-4}$ & 0.035829 & 0.015520 \\
4 & $4.656764\times10^{-7}$ & $1.298178\times10^{-7}$ & 0.047277 & 0.021596 \\
5 & $1.427445\times10^{-13}$ & $1.161032\times10^{-14}$ & 0.058077 & 0.026513 \\
6 & $1.489288\times10^{-26}$ & $9.508624\times10^{-29}$ & 0.068567 & 0.031778 \\
7 & $1.497717\times10^{-52}$ & $6.264169\times10^{-57}$ & 0.077042 & 0.037412 \\
8 & $1.607123\times10^{-104}$ & $2.755944\times10^{-113}$ & 0.086658 & 0.042781 \\
9 & $1.76902\times10^{-208}$ & $5.279411\times10^{-226}$ & 0.095481 & 0.048186 \\
10 & $2.217477\times10^{-416}$ & $1.952676\times10^{-451}$ & 0.102580 & 0.053230 \\
11 & $3.395322\times10^{-832}$ & $2.655362\times10^{-902}$ & 0.110673 & 0.058621 \\
\bottomrule
\end{tabular}
\caption{Product torus: complete trajectory of residuals and cumulative times. The times are cumulative medians of three alternating runs with 1200 decimal digits and tolerance $10^{-500}$.}
\label{tab:supp_product_torus}
\end{table}

\section{The Stiefel manifold: costly exponential map and simple retraction}
\label{subsec:stiefel}
% ------------------------------------------------------------

Consider the Stiefel manifold
\[
    \operatorname{St}(n,r)
    =
    \left\{
        Y\in\R^{n\times r}:
        Y^\top Y=I_r
    \right\}.
\]
This manifold is the regular level set of the mapping
\[
    F(Y)=Y^\top Y-I_r,
\]
understood as taking values in the space
$\operatorname{Sym}(r)$ of symmetric matrices. Its derivative is
\begin{equation}\label{eq:DF_stiefel}
    \D F(Y)(H)
    =
    Y^\top H+H^\top Y.
\end{equation}
This mapping is surjective. Indeed, if
$S\in\operatorname{Sym}(r)$, by taking
\[
    H=\frac12YS
\]
we obtain
\[
    \D F(Y)(H)
    =
    \frac12S+\frac12S
    =
    S.
\]
Therefore, $0$ is a regular value of $F$.

The tangent space is
\begin{equation}\label{eq:tangent_stiefel}
    \T_Y\operatorname{St}(n,r)
    =
    \left\{
        \eta\in\R^{n\times r}:
        Y^\top\eta+\eta^\top Y=0
    \right\}.
\end{equation}

In this subsection, we use the Euclidean metric induced by the
Frobenius inner product,
\[
    \langle U,V\rangle_F
    =
    \operatorname{tr}(U^\top V).
\]
The normal space is
\[
    N_Y\operatorname{St}(n,r)
    =
    \left\{
        YS:
        S\in\operatorname{Sym}(r)
    \right\},
\]
and the orthogonal projector onto the tangent space is
\begin{equation}\label{eq:projector_stiefel}
    \Proj_Y(Z)
    =
    Z-Y\,\operatorname{sym}(Y^\top Z),
\end{equation}
where
\[
    \operatorname{sym}(B)
    =
    \frac12(B+B^\top).
\]

% ------------------------------------------------------------
\subsection*{Riemannian exponential for the Euclidean metric}
% ------------------------------------------------------------

The Riemannian exponential of the induced Euclidean metric can be
computed by means of a block formula. Let
\[
    \eta\in \T_Y\operatorname{St}(n,r),
\]
and define
\[
    A=Y^\top\eta,
    \qquad
    S=\eta^\top\eta.
\]
Since $\eta$ is tangent,
\[
    A^\top=-A.
\]
The Euclidean geodesic with initial conditions
\[
    \gamma(0)=Y,
    \qquad
    \dot\gamma(0)=\eta
\]
satisfies
\begin{equation}\label{eq:euclidean_exp_stiefel}
\begin{aligned}
    \Exp^{E}_Y(\eta)
    &=
    \begin{bmatrix}
        Y&\eta
    \end{bmatrix}
    \exp
    \begin{pmatrix}
        A&-S\\
        I_r&A
    \end{pmatrix}
    \begin{pmatrix}
        \exp(-A)\\
        0
    \end{pmatrix}.
\end{aligned}
\end{equation}
This expression corresponds to the embedded, or Euclidean, metric and requires,
in general, a matrix exponential of size $2r\times2r$, in addition to
the exponential of the skew-symmetric matrix $A\in\R^{r\times r}$.

In the case $\operatorname{St}(3,2)$, formula
\eqref{eq:euclidean_exp_stiefel} requires computing one matrix exponential of
size $4\times4$ and another of size $2\times2$.

% ------------------------------------------------------------
\subsection*{Polar retraction}
% ------------------------------------------------------------

A much simpler alternative is the polar retraction
\begin{equation}\label{eq:polar_retraction_stiefel}
    R_Y^{\mathrm{pol}}(\eta)
    =
    (Y+\eta)
    \bigl(I_r+\eta^\top\eta\bigr)^{-1/2}.
\end{equation}
Indeed, using the tangency condition,
\[
\begin{aligned}
    (Y+\eta)^\top(Y+\eta)
    &=
    Y^\top Y
    +
    Y^\top\eta
    +
    \eta^\top Y
    +
    \eta^\top\eta\\
    &=
    I_r+\eta^\top\eta.
\end{aligned}
\]
Therefore,
\[
\begin{aligned}
    \bigl(R_Y^{\mathrm{pol}}(\eta)\bigr)^\top
    R_Y^{\mathrm{pol}}(\eta)
    &=
    \bigl(I_r+\eta^\top\eta\bigr)^{-1/2}
    \bigl(I_r+\eta^\top\eta\bigr)
    \bigl(I_r+\eta^\top\eta\bigr)^{-1/2}\\
    &=I_r.
\end{aligned}
\]
Thus,
\[
    R_Y^{\mathrm{pol}}(\eta)
    \in
    \operatorname{St}(n,r).
\]
Moreover,
\[
    R_Y^{\mathrm{pol}}(0)=Y,
    \qquad
    \D R_Y^{\mathrm{pol}}(0)(\eta)=\eta.
\]
Therefore, it is a retraction and, in particular, it satisfies the first-order
assumptions used in the convergence theorems.

Its geometric cost is concentrated in the inverse square root of an
$r\times r$ matrix. For $\operatorname{St}(3,2)$, only the inverse square root
of a symmetric positive-definite matrix of size $2\times2$ is required.

% ------------------------------------------------------------
\subsection*{A nonlinear and coupled field on $\operatorname{St}(3,2)$}
% ------------------------------------------------------------

From now on, consider
\[
    Y=
    \begin{pmatrix}
        y_{11}&y_{12}\\
        y_{21}&y_{22}\\
        y_{31}&y_{32}
    \end{pmatrix}
    \in\operatorname{St}(3,2).
\]
We first define three linear forms:
\begin{align}
    u_1(Y)
    &=
    y_{11}
    +\frac25y_{22}
    -\frac3{10}y_{31}
    +\frac15y_{12},
    \label{eq:u1_stiefel}\\
    u_2(Y)
    &=
    y_{21}
    -\frac15y_{12}
    +\frac12y_{32}
    +\frac1{10}y_{11},
    \label{eq:u2_stiefel}\\
    u_3(Y)
    &=
    y_{31}
    +\frac3{10}y_{22}
    -\frac25y_{21}
    +\frac15y_{32}.
    \label{eq:u3_stiefel}
\end{align}
From them, we construct the nonlinear functions
\begin{align}
    g_1(Y)
    &=
    u_1+u_2^2-1.6119,
    \label{eq:g1_stiefel}\\
    g_2(Y)
    &=
    u_2+u_1u_3-1.2468,
    \label{eq:g2_stiefel}\\
    g_3(Y)
    &=
    u_3+u_1^2+\frac12u_2u_3-1.6423.
    \label{eq:g3_stiefel}
\end{align}
These expressions depend on the six entries of $Y$ and are coupled.
The decimal constants are fixed as part of the definition of the problem;
they are not adjusted after running the method. Their choice produces a
nontrivial zero, away from the coordinate frames, and a regular reduced matrix.

For
\[
    v=
    \begin{pmatrix}
        v_1\\v_2\\v_3
    \end{pmatrix},
\]
denote by
\[
    \widehat v
    =
    \begin{pmatrix}
        0&-v_3&v_2\\
        v_3&0&-v_1\\
        -v_2&v_1&0
    \end{pmatrix}
\]
the associated skew-symmetric matrix. We define
\begin{equation}\label{eq:field_stiefel}
    X(Y)
    =
    \widehat{g(Y)}\,Y,
    \qquad
    g(Y)=
    \begin{pmatrix}
        g_1(Y)\\g_2(Y)\\g_3(Y)
    \end{pmatrix}.
\end{equation}

The field is tangent. Indeed,
\[
\begin{aligned}
    Y^\top X(Y)+X(Y)^\top Y
    &=
    Y^\top\widehat{g(Y)}Y
    +
    Y^\top\widehat{g(Y)}^\top Y\\
    &=
    Y^\top
    \left(
        \widehat{g(Y)}
        -
        \widehat{g(Y)}
    \right)
    Y\\
    &=0.
\end{aligned}
\]

Moreover,
\[
    X(Y)=0
    \quad\Longleftrightarrow\quad
    g(Y)=0.
\]
Indeed, if $\widehat{g(Y)}\neq0$, its null space in $\R^3$ has dimension
one and cannot contain the two linearly independent columns of
$Y\in\operatorname{St}(3,2)$.

% ------------------------------------------------------------
\subsection*{Ambient derivative and covariant derivative}
% ------------------------------------------------------------

Let
\[
    H=
    \begin{pmatrix}
        h_{11}&h_{12}\\
        h_{21}&h_{22}\\
        h_{31}&h_{32}
    \end{pmatrix}.
\]
The derivatives of the linear forms are
\begin{align*}
    \D u_1(Y)(H)
    &=
    h_{11}
    +\frac25h_{22}
    -\frac3{10}h_{31}
    +\frac15h_{12},\\
    \D u_2(Y)(H)
    &=
    h_{21}
    -\frac15h_{12}
    +\frac12h_{32}
    +\frac1{10}h_{11},\\
    \D u_3(Y)(H)
    &=
    h_{31}
    +\frac3{10}h_{22}
    -\frac25h_{21}
    +\frac15h_{32}.
\end{align*}
If we write
\[
    \D u(Y)(H)
    =
    \begin{pmatrix}
        \delta u_1\\
        \delta u_2\\
        \delta u_3
    \end{pmatrix},
\]
then
\begin{equation}\label{eq:Dg_stiefel}
    \D g(Y)(H)
    =
    \begin{pmatrix}
        1&2u_2&0\\
        u_3&1&u_1\\
        2u_1&\frac12u_3&1+\frac12u_2
    \end{pmatrix}
    \D u(Y)(H).
\end{equation}
By the product rule,
\begin{equation}\label{eq:ambient_DX_stiefel}
    \D\widetilde X(Y)(H)
    =
    \widehat{\D g(Y)(H)}\,Y
    +
    \widehat{g(Y)}\,H.
\end{equation}
The covariant derivative for the induced Euclidean metric is
\begin{equation}\label{eq:covariant_DX_stiefel}
    \nabla_H X(Y)
    =
    \Proj_Y\left(
        \widehat{\D g(Y)(H)}\,Y
        +
        \widehat{g(Y)}\,H
    \right).
\end{equation}

% ------------------------------------------------------------
\subsection*{Reduced $3\times3$ system}
% ------------------------------------------------------------

Consider the matrices
\[
    K_1=
    \begin{pmatrix}
        0&0&0\\
        0&0&-1\\
        0&1&0
    \end{pmatrix},
    \qquad
    K_2=
    \begin{pmatrix}
        0&0&1\\
        0&0&0\\
        -1&0&0
    \end{pmatrix},
\]
\[
    K_3=
    \begin{pmatrix}
        0&-1&0\\
        1&0&0\\
        0&0&0
    \end{pmatrix}.
\]
We define
\begin{equation}\label{eq:tangent_basis_stiefel}
    B_i(Y)=K_iY,
    \qquad
    i=1,2,3.
\end{equation}
Since $\operatorname{St}(3,2)$ has dimension three, these vectors form
a global basis of $\T_Y\operatorname{St}(3,2)$.

We write
\[
    \eta
    =
    \xi_1B_1(Y)
    +
    \xi_2B_2(Y)
    +
    \xi_3B_3(Y).
\]
The Newton equation reduces to the system
\begin{equation}\label{eq:system_3x3_stiefel}
    M(Y)
    \begin{pmatrix}
        \xi_1\\
        \xi_2\\
        \xi_3
    \end{pmatrix}
    =
    -b(Y),
\end{equation}
where
\[
    M_{ij}(Y)
    =
    \left\langle
        B_i(Y),
        \nabla_{B_j(Y)}X(Y)
    \right\rangle_F,
\]
and
\[
    b_i(Y)
    =
    \langle B_i(Y),X(Y)\rangle_F.
\]
Once \eqref{eq:system_3x3_stiefel} has been solved,
\[
    \eta
    =
    \sum_{j=1}^3\xi_jB_j(Y).
\]

% ------------------------------------------------------------

\subsection*{Updated numerical results}
\begin{table}[H]
\centering
\scriptsize
\setlength{\tabcolsep}{4.2pt}
\begin{tabular}{rccrr}
\toprule
$k$ & $\|X(p_k)\|_{\Exp}$ & $\|X(p_k)\|_R$
& $t_{\Exp,k}$ (s) & $t_{R,k}$ (s)\\
\midrule
0 & $9.758255\times10^{-1}$ & $9.758255\times10^{-1}$ & 0.001169 & 0.001196 \\
1 & $8.201328\times10^{-1}$ & $6.478493\times10^{-1}$ & 0.088819 & 0.024914 \\
2 & $1.077263\times10^{-1}$ & $1.037519\times10^{-1}$ & 0.165888 & 0.047424 \\
3 & $2.286715\times10^{-2}$ & $1.882023\times10^{-2}$ & 0.233804 & 0.072950 \\
4 & $1.184656\times10^{-3}$ & $7.663722\times10^{-4}$ & 0.297837 & 0.094344 \\
5 & $3.656274\times10^{-6}$ & $1.572548\times10^{-6}$ & 0.358700 & 0.115600 \\
6 & $3.899457\times10^{-11}$ & $7.277254\times10^{-12}$ & 0.407913 & 0.137056 \\
7 & $4.471085\times10^{-21}$ & $1.555223\times10^{-22}$ & 0.449970 & 0.158501 \\
8 & $5.845408\times10^{-41}$ & $7.069257\times10^{-44}$ & 0.486054 & 0.180165 \\
9 & $9.986523\times10^{-81}$ & $1.460687\times10^{-86}$ & 0.520960 & 0.207514 \\
10 & $2.915642\times10^{-160}$ & $6.237787\times10^{-172}$ & 0.556403 & 0.229489 \\
11 & $2.485343\times10^{-319}$ & $1.13757\times10^{-342}$ & 0.592360 & 0.250823 \\
12 & $1.805863\times10^{-637}$ & $3.783271\times10^{-684}$ & 0.626350 & 0.271646 \\
\bottomrule
\end{tabular}
\caption{Nonlinear Stiefel problem: complete trajectory of residuals and cumulative times. The times are cumulative medians of three alternating runs with 1200 decimal digits and tolerance $10^{-500}$.}
\label{tab:supp_stiefel_nonlinear}
\end{table}

\section{Complete multiprecision trajectories}
\label{sec:supp_trajectories}

The multiprecision protocol uses 1200 decimal digits, tolerance
$10^{-500}$, and three alternating repetitions of each variant. In each row,
the residual and cumulative time belong to the same protocol. The times
are the cumulative medians by iteration. If one variant terminates earlier,
no subsequent iterations are invented; for this reason, dashes appear for the
retraction in the example on $S^3$ after its one-step termination.

\begin{table}[H]
\centering
\scriptsize
\setlength{\tabcolsep}{4.2pt}
\begin{tabular}{rccrr}
\toprule
$k$ & $\|X(p_k)\|_{\Exp}$ & $\|X(p_k)\|_R$
& $t_{\Exp,k}$ (s) & $t_{R,k}$ (s)\\
\midrule
0 & $7.874643\times10^{-1}$ & $7.874643\times10^{-1}$ & 0.000268 & 0.000256 \\
1 & $1.434361\times10^{-1}$ & $1.436173\times10^{-1}$ & 0.006942 & 0.003726 \\
2 & $1.267638\times10^{-2}$ & $1.267914\times10^{-2}$ & 0.013534 & 0.006721 \\
3 & $1.524631\times10^{-4}$ & $1.524634\times10^{-4}$ & 0.019511 & 0.009878 \\
4 & $2.323057\times10^{-8}$ & $2.323057\times10^{-8}$ & 0.025137 & 0.013068 \\
5 & $5.396595\times10^{-16}$ & $5.396595\times10^{-16}$ & 0.030297 & 0.016241 \\
6 & $2.912324\times10^{-31}$ & $2.912324\times10^{-31}$ & 0.034642 & 0.019385 \\
7 & $8.481631\times10^{-62}$ & $8.481631\times10^{-62}$ & 0.038993 & 0.022638 \\
8 & $7.193807\times10^{-123}$ & $7.193807\times10^{-123}$ & 0.043338 & 0.025694 \\
9 & $5.175086\times10^{-245}$ & $5.175086\times10^{-245}$ & 0.047615 & 0.028676 \\
10 & $2.678152\times10^{-489}$ & $2.678152\times10^{-489}$ & 0.051941 & 0.031653 \\
11 & $7.172496\times10^{-978}$ & $7.172496\times10^{-978}$ & 0.057418 & 0.034571 \\
\bottomrule
\end{tabular}
\caption{Cylinder: complete trajectory of residuals and cumulative times. The times are cumulative medians of three alternating runs with 1200 decimal digits and tolerance $10^{-500}$.}
\label{tab:supp_cylinder}
\end{table}

\begin{table}[H]
\centering
\scriptsize
\setlength{\tabcolsep}{4.2pt}
\begin{tabular}{rccrr}
\toprule
$k$ & $\|X(p_k)\|_{\Exp}$ & $\|X(p_k)\|_R$
& $t_{\Exp,k}$ (s) & $t_{R,k}$ (s)\\
\midrule
0 & $5.0\times10^{-1}$ & $5.0\times10^{-1}$ & 0.001096 & 0.000922 \\
1 & $5.372561\times10^{-2}$ & $3.363705\times10^{-1202}$ & 0.016217 & 0.014422 \\
2 & $5.182658\times10^{-5}$ & -- & 0.029973 & -- \\
3 & $4.640197\times10^{-14}$ & -- & 0.043548 & -- \\
4 & $3.330336\times10^{-41}$ & -- & 0.057234 & -- \\
5 & $1.23124\times10^{-122}$ & -- & 0.071222 & -- \\
6 & $6.221675\times10^{-367}$ & -- & 0.084759 & -- \\
7 & $8.027875\times10^{-1100}$ & -- & 0.098664 & -- \\
\bottomrule
\end{tabular}
\caption{$S^3$: complete trajectory of residuals and cumulative times. The times are cumulative medians of three alternating runs with 1200 decimal digits and tolerance $10^{-500}$.}
\label{tab:supp_s3}
\end{table}

\begin{table}[H]
\centering
\scriptsize
\setlength{\tabcolsep}{4.2pt}
\begin{tabular}{rccrr}
\toprule
$k$ & $\|X(p_k)\|_{\Exp}$ & $\|X(p_k)\|_R$
& $t_{\Exp,k}$ (s) & $t_{R,k}$ (s)\\
\midrule
0 & $5.484387\times10^{-1}$ & $5.484387\times10^{-1}$ & 0.000672 & 0.000626 \\
1 & $3.607681\times10^{-1}$ & $3.256465\times10^{-1}$ & 2.021257 & 0.016090 \\
2 & $2.326791\times10^{-2}$ & $3.606411\times10^{-2}$ & 2.977476 & 0.029715 \\
3 & $4.350643\times10^{-4}$ & $3.564352\times10^{-4}$ & 3.207198 & 0.042812 \\
4 & $3.951101\times10^{-8}$ & $2.632585\times10^{-8}$ & 3.399987 & 0.056178 \\
5 & $4.043158\times10^{-16}$ & $1.476852\times10^{-16}$ & 3.598228 & 0.069326 \\
6 & $3.034939\times10^{-32}$ & $4.713945\times10^{-33}$ & 3.790723 & 0.082565 \\
7 & $1.551148\times10^{-64}$ & $4.778443\times10^{-66}$ & 4.006702 & 0.096186 \\
8 & $4.170436\times10^{-129}$ & $4.861653\times10^{-132}$ & 4.198194 & 0.110084 \\
9 & $3.053869\times10^{-258}$ & $5.01734\times10^{-264}$ & 4.394884 & 0.124258 \\
10 & $1.636597\times10^{-516}$ & $5.364794\times10^{-528}$ & 4.586303 & 0.137583 \\
\bottomrule
\end{tabular}
\caption{$SL(2,\R)$: Riemannian exponential and homogeneous retraction. The times are cumulative medians of three alternating runs with 1200 decimal digits and tolerance $10^{-500}$.}
\label{tab:supp_sl2}
\end{table}

\begin{table}[H]
\centering
\scriptsize
\setlength{\tabcolsep}{4.2pt}
\begin{tabular}{rccrr}
\toprule
$k$ & $\|X(p_k)\|_{\Exp}$ & $\|X(p_k)\|_R$
& $t_{\Exp,k}$ (s) & $t_{R,k}$ (s)\\
\midrule
0 & $2.606035\times10^{-1}$ & $2.606035\times10^{-1}$ & 0.001357 & 0.001476 \\
1 & $1.762323\times10^{-2}$ & $1.28056\times10^{-2}$ & 0.016806 & 0.013220 \\
2 & $5.149984\times10^{-6}$ & $1.537462\times10^{-6}$ & 0.031659 & 0.025334 \\
3 & $1.419023\times10^{-16}$ & $3.006613\times10^{-18}$ & 0.045072 & 0.037162 \\
4 & $3.218826\times10^{-48}$ & $2.446631\times10^{-53}$ & 0.058293 & 0.049224 \\
5 & $3.986347\times10^{-143}$ & $1.38595\times10^{-158}$ & 0.071708 & 0.061453 \\
6 & $7.883916\times10^{-428}$ & $2.588342\times10^{-474}$ & 0.084551 & 0.073689 \\
7 & $4.119681\times10^{-1201}$ & $2.464168\times10^{-1201}$ & 0.098081 & 0.085713 \\
\bottomrule
\end{tabular}
\caption{Wahba problem: complete trajectory of residuals and cumulative times. The times are cumulative medians of three alternating runs with 1200 decimal digits and tolerance $10^{-500}$.}
\label{tab:supp_wahba}
\end{table}

\begin{table}[H]
\centering
\scriptsize
\setlength{\tabcolsep}{4.2pt}
\begin{tabular}{rccrr}
\toprule
$k$ & $\|X(p_k)\|_{\Exp}$ & $\|X(p_k)\|_R$
& $t_{\Exp,k}$ (s) & $t_{R,k}$ (s)\\
\midrule
0 & $2.339125\times10^{0}$ & $2.339125\times10^{0}$ & 0.004200 & 0.003970 \\
1 & $3.810634\times10^{-1}$ & $1.492113\times10^{-1}$ & 0.036781 & 0.030213 \\
2 & $1.49713\times10^{-3}$ & $2.94527\times10^{-5}$ & 0.068347 & 0.060097 \\
3 & $1.441875\times10^{-10}$ & $2.597738\times10^{-16}$ & 0.099895 & 0.085820 \\
4 & $1.307331\times10^{-31}$ & $1.550253\times10^{-48}$ & 0.128240 & 0.110914 \\
5 & $9.642509\times10^{-95}$ & $5.45052\times10^{-145}$ & 0.156681 & 0.136953 \\
6 & $3.898301\times10^{-284}$ & $2.61254\times10^{-434}$ & 0.186891 & 0.163861 \\
7 & $2.56216\times10^{-852}$ & $1.254982\times10^{-1201}$ & 0.215074 & 0.188965 \\
\bottomrule
\end{tabular}
\caption{Toroidal von Mises model: complete trajectory of residuals and cumulative times. The times are cumulative medians of three alternating runs with 1200 decimal digits and tolerance $10^{-500}$.}
\label{tab:supp_von_mises}
\end{table}

\begin{table}[H]
\centering
\scriptsize
\setlength{\tabcolsep}{4.2pt}
\begin{tabular}{rccrr}
\toprule
$k$ & $\|X(p_k)\|_{\Exp}$ & $\|X(p_k)\|_R$
& $t_{\Exp,k}$ (s) & $t_{R,k}$ (s)\\
\midrule
0 & $4.719654\times10^{-1}$ & $4.719654\times10^{-1}$ & 0.002236 & 0.002561 \\
1 & $5.353999\times10^{-2}$ & $4.625133\times10^{-2}$ & 0.086563 & 0.032679 \\
2 & $5.624617\times10^{-4}$ & $4.930562\times10^{-4}$ & 0.160845 & 0.063338 \\
3 & $4.173036\times10^{-8}$ & $2.588569\times10^{-8}$ & 0.222026 & 0.088974 \\
4 & $3.18385\times10^{-16}$ & $1.457545\times10^{-16}$ & 0.269569 & 0.115542 \\
5 & $1.33614\times10^{-32}$ & $2.260606\times10^{-33}$ & 0.317738 & 0.144748 \\
6 & $3.264013\times10^{-65}$ & $1.11161\times10^{-66}$ & 0.357374 & 0.173028 \\
7 & $1.404269\times10^{-130}$ & $1.314879\times10^{-133}$ & 0.398799 & 0.200872 \\
8 & $3.605363\times10^{-261}$ & $3.760746\times10^{-267}$ & 0.439157 & 0.230718 \\
9 & $1.713344\times10^{-522}$ & $1.504975\times10^{-534}$ & 0.479302 & 0.256425 \\
\bottomrule
\end{tabular}
\caption{Brockett eigenproblem: complete trajectory of residuals and cumulative times. The times are cumulative medians of three alternating runs with 1200 decimal digits and tolerance $10^{-500}$.}
\label{tab:supp_brockett}
\end{table}

\section{Timing summary and moderate-precision protocol}
\label{sec:supp_timing_summary}

\begin{table}[H]
\centering
\scriptsize
\setlength{\tabcolsep}{3.2pt}
\resizebox{\textwidth}{!}{%
\begin{tabular}{lrrrrrrr}
\toprule
Problem & $k_\Exp$ & $k_R$ & $t_\Exp$ & MAD$_\Exp$ & $t_R$ & MAD$_R$ & Reduction\\
\midrule
Cylinder & 11 & 11 & 0.057418 & 0.000068 & 0.034571 & 0.000360 & 39.8\% \\
Sphere $S^3$ & 7 & 1 & 0.098664 & 0.000530 & 0.014422 & 0.000166 & 85.4\% \\
Product torus & 11 & 11 & 0.110673 & 0.003233 & 0.058621 & 0.001686 & 47.0\% \\
Nonlinear $SL(2,\R)$ & 10 & 10 & 4.586303 & 0.093026 & 0.137583 & 0.003157 & 97.0\% \\
Nonlinear Stiefel problem & 12 & 12 & 0.626350 & 0.003254 & 0.271646 & 0.006032 & 56.6\% \\
Wahba on $S^3$ & 7 & 7 & 0.098081 & 0.003024 & 0.085713 & 0.001819 & 12.6\% \\
Toroidal von Mises & 7 & 7 & 0.215074 & 0.003186 & 0.188965 & 0.011475 & 12.1\% \\
Brockett on $\operatorname{St}(3,2)$ & 9 & 9 & 0.479302 & 0.013373 & 0.256425 & 0.010924 & 46.5\% \\
\bottomrule
\end{tabular}}
\caption{Summary of the multiprecision protocol. Times are expressed in seconds.}
\label{tab:supp_multiprecision_summary}
\end{table}

\Cref{tab:supp_multiprecision_summary} answers the question of how much it costs to
reach the extreme tolerance used to study the asymptotic regime.
It should not be interpreted as a universal benchmark: the cost of transcendental
functions, matrix exponentials, and geodesic integration increases with the
requested precision.

As an additional control, a second protocol was run with 50 digits and
tolerance $10^{-14}$. Each method was warmed up twice and measured over
eleven repetitions, alternating the order of execution. The median and the
median absolute deviation (MAD) of the total time are reported.

\begin{table}[H]
\centering
\scriptsize
\setlength{\tabcolsep}{3.2pt}
\resizebox{\textwidth}{!}{%
\begin{tabular}{lrrrrrrr}
\toprule
Problem & $k_\Exp$ & $k_R$ & $t_\Exp$ & MAD$_\Exp$ & $t_R$ & MAD$_R$ & Reduction\\
\midrule
Cylinder & 5 & 5 & 0.004506 & 0.000485 & 0.003922 & 0.000095 & 13.0\% \\
Sphere $S^3$ & 4 & 1 & 0.010709 & 0.000462 & 0.002558 & 0.000058 & 76.1\% \\
Product torus & 6 & 6 & 0.009887 & 0.000843 & 0.008421 & 0.000281 & 14.8\% \\
Nonlinear $SL(2,\R)$ & 5 & 5 & 0.172084 & 0.005870 & 0.015709 & 0.000640 & 90.9\% \\
Nonlinear Stiefel problem & 7 & 7 & 0.057641 & 0.000564 & 0.033332 & 0.000638 & 42.2\% \\
Wahba on $S^3$ & 3 & 3 & 0.008871 & 0.000848 & 0.008168 & 0.000637 & 7.9\% \\
Toroidal von Mises & 4 & 3 & 0.005742 & 0.000076 & 0.004185 & 0.000081 & 27.1\% \\
Brockett on $\operatorname{St}(3,2)$ & 4 & 4 & 0.042311 & 0.001870 & 0.027628 & 0.001075 & 34.7\% \\
\bottomrule
\end{tabular}}
\caption{Total times at moderate precision: 50 digits, tolerance $10^{-14}$,
two warm-up runs, and eleven alternating repetitions.}
\label{tab:supp_moderate_summary}
\end{table}

The clearest reductions in both protocols appear for $SL(2,\R)$ and in
the Stiefel problems. In the low-cost experiments, the differences should be
interpreted together with the MAD and not as a universal superiority.

\section[Controls for the Riemannian exponential on SL(2,R)]{Controls for the Riemannian exponential on $SL(2,\R)$}
\label{sec:supp_sl2_diagnostics}

The Riemannian exponential is evaluated by integrating the geodesic equation of
the induced Frobenius metric. The implementation uses a recurrent Taylor series
and adapts the tolerance to the residual of the Newton step. For each geodesic,
the truncation order, an a posteriori estimate based on the recent ratios of
the coefficients, and controls of feasibility, tangency, the geodesic equation,
and energy conservation are recorded.

\begin{table}[H]
\centering
\scriptsize
\setlength{\tabcolsep}{3.5pt}
\resizebox{\textwidth}{!}{%
\begin{tabular}{rccrcc}
\toprule
Step & Residual before the step & Requested tolerance & Series order
& Tail estimate & Recent ratio\\
\midrule
0 & $5.4844\times10^{-1}$ & $1.0\times10^{-80}$ & 109 & $9.007\times10^{-81}$ & $1.96\times10^{-1}$ \\
1 & $3.6077\times10^{-1}$ & $1.0\times10^{-80}$ & 76 & $3.078\times10^{-81}$ & $9.795\times10^{-2}$ \\
2 & $2.3268\times10^{-2}$ & $1.0\times10^{-80}$ & 36 & $6.739\times10^{-81}$ & $6.723\times10^{-3}$ \\
3 & $4.3506\times10^{-4}$ & $1.0\times10^{-80}$ & 33 & $2.855\times10^{-136}$ & $9.931\times10^{-5}$ \\
4 & $3.9511\times10^{-8}$ & $1.0\times10^{-80}$ & 33 & $2.528\times10^{-271}$ & $7.773\times10^{-9}$ \\
5 & $4.0432\times10^{-16}$ & $1.0\times10^{-80}$ & 33 & $1.492\times10^{-536}$ & $7.196\times10^{-17}$ \\
6 & $3.0349\times10^{-32}$ & $2.795\times10^{-95}$ & 33 & $7.279\times10^{-1067}$ & $5.971\times10^{-33}$ \\
7 & $1.5511\times10^{-64}$ & $3.732\times10^{-192}$ & 33 & $1.081\times10^{-2131}$ & $3.324\times10^{-65}$ \\
8 & $4.1704\times10^{-129}$ & $7.253\times10^{-386}$ & 33 & $9.368\times10^{-4263}$ & $8.618\times10^{-130}$ \\
9 & $3.0539\times10^{-258}$ & $1.0\times10^{-560}$ & 33 & $2.075\times10^{-8524}$ & $6.373\times10^{-259}$ \\
\bottomrule
\end{tabular}}
\caption{Data for the a posteriori truncation criterion for the geodesics of
$SL(2,\R)$. The tail estimate is not presented as a rigorous certification
of the entire future series.}
\label{tab:supp_sl2_truncation}
\end{table}

\begin{table}[H]
\centering
\scriptsize
\setlength{\tabcolsep}{3.5pt}
\resizebox{\textwidth}{!}{%
\begin{tabular}{rccccc}
\toprule
Step & Raw $\det$ defect & Stabilization & Tangency defect
& Geodesic residual & Energy defect\\
\midrule
0 & $5.979\times10^{-82}$ & $4.997\times10^{-82}$ & $6.568\times10^{-80}$ & $1.687\times10^{-77}$ & $6.34\times10^{-80}$ \\
1 & $2.843\times10^{-82}$ & $2.354\times10^{-82}$ & $2.188\times10^{-80}$ & $1.54\times10^{-78}$ & $2.773\times10^{-81}$ \\
2 & $5.1\times10^{-83}$ & $4.217\times10^{-83}$ & $1.887\times10^{-81}$ & $5.889\times10^{-80}$ & $5.844\times10^{-84}$ \\
3 & $4.275\times10^{-141}$ & $3.535\times10^{-141}$ & $1.454\times10^{-139}$ & $2.686\times10^{-137}$ & $1.624\times10^{-142}$ \\
4 & $2.85\times10^{-279}$ & $2.357\times10^{-279}$ & $9.691\times10^{-278}$ & $2.155\times10^{-276}$ & $2.088\times10^{-286}$ \\
5 & $1.663\times10^{-552}$ & $1.375\times10^{-552}$ & $5.655\times10^{-551}$ & $1.194\times10^{-549}$ & $1.187\times10^{-567}$ \\
6 & $6.523\times10^{-1099}$ & $5.393\times10^{-1099}$ & $2.218\times10^{-1097}$ & $4.927\times10^{-1096}$ & $6.718\times10^{-1130}$ \\
7 & $1.554\times10^{-1201}$ & $0$ & $8.852\times10^{-1266}$ & $1.096\times10^{-1330}$ & $0$ \\
8 & $7.768\times10^{-1202}$ & $1.79\times10^{-1201}$ & $1.681\times10^{-1330}$ & $4.526\times10^{-1460}$ & $8.084\times10^{-1460}$ \\
9 & $1.554\times10^{-1201}$ & $0$ & $8.084\times10^{-1460}$ & $4.835\times10^{-1718}$ & $8.412\times10^{-1718}$ \\
\bottomrule
\end{tabular}}
\caption{Geometric and dynamical controls for the integration of the
Riemannian exponential on $SL(2,\R)$.}
\label{tab:supp_sl2_controls}
\end{table}

The final normalization of the determinant is used only to correct the
rounding error of the integration. Its norm is recorded in
\Cref{tab:supp_sl2_controls}. The code does not use the group exponential.

\section{Implementation and reproducibility}
\label{sec:supp_reproducibility}
\begingroup\small

The repository contains eight standalone experiment programs and the global
result generator \path{src/generate_article_tables.py}. The generator accepts
the option \path{--output-root} and creates, under the selected destination,
the directories \path{multiprecision}, \path{moderate_precision}, and
\path{diagnostics}. It also writes the file
\path{execution_environment.json}.

The reference execution archived in the repository was carried out with:
\begin{itemize}[nosep,leftmargin=*]
    \item Python: \texttt{3.13.5};
    \item \texttt{mpmath}: \texttt{1.3.0};
    \item platform: \nolinkurl{Linux-4.4.0-x86_64-with-glibc2.41};
    \item timer: \nolinkurl{time.perf_counter};
    \item UTC date: \nolinkurl{2026-07-16T00:18:43.319220+00:00}.
\end{itemize}

From the repository root, all numerical results are regenerated by
\begin{verbatim}
python src/generate_article_tables.py --output-root results_local
\end{verbatim}
This command writes a local regeneration to \path{results_local} and leaves
the committed reference data in \path{results} unchanged. The exact numerical
dependency is pinned in \path{requirements.txt}. Detailed execution instructions
are provided in \path{docs/REPRODUCIBILITY.md}, and the two numerical protocols
are documented in \path{docs/NUMERICAL_PROTOCOLS.md}. The committed reference
files are stored in \path{results/multiprecision},
\path{results/moderate_precision}, and \path{results/diagnostics}; the runtime
metadata are stored in \path{results/execution_environment.json}. Repository
integrity can be checked with \path{scripts/verify_repository.py}, while the
SHA-256 manifest is available at \path{checksums/SHA256SUMS.txt}.

The programs verify, depending on the problem, tangency of the direction,
the residual of the reduced system, satisfaction of the constraints, coincidence
of the final roots, and regularity of the zero. No number in the manuscript or
Online Resource 1 is adjusted manually.
\endgroup
\end{document}
